\chapter{Production DataOps and Observability: Terraform, dbt, Schemachange, and Alerts}
\index{DataOps}\index{Terraform}\index{dbt}\index{schemachange}\index{Snowflake Alerts}\index{Week 23}%
\begin{objectives}
\begin{itemize}
    \item Provision Snowflake infrastructure as code with Terraform and manage credit spend with resource monitors.
    \item Enforce data quality with dbt generic and singular tests plus Snowflake Data Metric Functions (DMFs).
    \item Run versioned, backward-compatible migrations with schemachange and the expand--migrate--contract pattern.
    \item Build Snowflake Alerts for cost anomalies over \texttt{ACCOUNT\_USAGE} and wire them to notification integrations.
    \item Assemble a hands-on GitHub Actions CI/CD pipeline deploying changes and testing models on pull requests.
    \item Architect a multi-environment (Dev/QA/Prod) deployment framework with zero-downtime schema evolution.
\end{itemize}
\end{objectives}

\begin{crosscloud}
DataOps is the most cloud-portable chapter in this book: Terraform, dbt, and GitHub Actions are the same tools on every platform; only the provider plugin and the warehouse-specific monitors differ.

\vspace{2mm}
{\footnotesize
\begin{tabular}{@{}p{2.8cm}p{3.0cm}p{3.0cm}p{3.0cm}@{}}
\toprule
\textbf{Capability} & \textbf{AWS} & \textbf{Azure} & \textbf{GCP} \\
\midrule
Infrastructure as code & Terraform AWS provider / CloudFormation & Terraform azurerm / Bicep & Terraform google provider \\
Warehouse IaC objects & Redshift clusters, S3, IAM & Synapse/Fabric workspaces & BigQuery datasets, slots \\
Transformation + tests & dbt-redshift & dbt-fabric / dbt-synapse & dbt-bigquery \\
Migration runner & Flyway / Liquibase / schemachange-style & Flyway / DbUp & Flyway / Liquibase \\
Spend monitoring & Cost Explorer + budgets & Cost Management + alerts & Budgets + billing export \\
\addlinespace
Snowflake equivalent & \multicolumn{3}{l}{terraform-provider-snowflake, dbt-snowflake, schemachange, resource monitors, alerts, DMFs} \\
\bottomrule
\end{tabular}}

\vspace{1mm}
\textbf{MongoDB} uses the same Terraform/dbt-adjacent discipline with Atlas providers; \textbf{Fabric} adds workspace deployment pipelines. The Snowflake-specific craft is in the platform-native observability: resource monitors, alerts, DMFs, and the \texttt{ACCOUNT\_USAGE} schema as the telemetry source of truth.
\end{crosscloud}

\section{Week 23 Roadmap}
Twenty-two chapters built capability; this one builds the \emph{operating system} around it. DataOps applies software engineering discipline---version control, code review, automated testing, CI/CD, observability---to data platform change, so that a Tuesday afternoon deploy is boring \cite{terraformDocs,dbtCoreDocs}.

Monday covers Terraform-managed infrastructure and resource monitors. Tuesday is dbt testing and DMFs. Wednesday handles migrations, alerts, and cost-anomaly detection. Thursday assembles the CI/CD pipeline. Friday designs the multi-environment framework.

\begin{keyterms}
\textbf{DataOps}: applying DevOps discipline---versioning, testing, CI/CD, observability---to data platform change.\quad
\textbf{Resource monitor}: a Snowflake object capping warehouse credit consumption with notify/suspend triggers.\quad
\textbf{DMF}: a Data Metric Function---a scheduled, table-attached data-quality metric (freshness, nulls, volume).\quad
\textbf{schemachange}: a versioned migration runner applying SQL change scripts in order, exactly once.\quad
\textbf{Expand--contract}: zero-downtime schema evolution: add new shape, migrate consumers, remove old shape---in separate releases.
\end{keyterms}

\section{Monday: Terraform-Managed Infrastructure and Resource Monitors}
\index{Terraform}\index{resource monitor}
Click-ops does not scale and cannot be reviewed. Terraform's Snowflake provider makes roles, warehouses, databases, monitors, and grants declarative: the desired state lives in git, \texttt{terraform plan} is the code review, and \texttt{apply} is the audited change \cite{terraformSnowflakeProvider}.

\begin{lstlisting}[caption={Terraform: warehouse plus a resource monitor with two triggers}]
resource "snowflake_warehouse" "etl_wh" {
  name           = "ETL_WH"
  warehouse_size = "XSMALL"
  auto_suspend   = 120
  auto_resume    = true
}

resource "snowflake_resource_monitor" "etl_budget" {
  name         = "ETL_MONTHLY_BUDGET"
  credit_quota = 200
  frequency    = "MONTHLY"
  triggers {
    on_percent = 75
    do_action  = "NOTIFY"
  }
  triggers {
    on_percent = 100
    do_action  = "SUSPEND_IMMEDIATE"
  }
  warehouses = [snowflake_warehouse.etl_wh.name]
}
\end{lstlisting}

The \textbf{resource monitor} is the platform-native cost circuit breaker: a credit quota with trigger percentages mapped to actions---notify, suspend (after running queries finish), or suspend immediately \cite{snowflakeResourceMonitorsDocs}. Every warehouse in a governed account has one; the 75\%-notify / 100\%-suspend-immediately pair is the standard posture for non-production, with production monitors tuned to page rather than halt.

\begin{pitfall}
\textbf{A suspended production warehouse is a self-inflicted outage.} Production monitors notify early and suspend late (or never, paging instead); development monitors suspend aggressively. Encode the difference in Terraform variables per environment so it cannot drift.
\end{pitfall}

\section{Tuesday: dbt Testing and Data Metric Functions}
\index{dbt}\index{DMF}\index{data quality}
dbt brings software engineering to SQL: models are versioned SELECTs with a dependency graph, and \emph{tests} make data quality executable. The four generic tests every model should carry---\texttt{unique}, \texttt{not\_null}, \texttt{accepted\_values}, \texttt{relationships}---plus singular tests for business invariants:

\begin{lstlisting}[language=SQL, caption={dbt incremental model plus schema.yml tests}]
-- models/marts/fct_orders.sql
{{ config(materialized='incremental', unique_key='order_id') }}
SELECT order_id, customer_id, order_ts, order_total, status
FROM {{ ref('stg_orders') }}
{% if is_incremental() %}
WHERE order_ts > (SELECT MAX(order_ts) FROM {{ this }})
{% endif %}

-- models/marts/schema.yml
-- version: 2
-- models:
--   - name: fct_orders
--     tests:
--       - dbt_utils.unique_combination_of_columns:
--           combination_of_columns: [order_id]
--     columns:
--       - name: order_id
--         tests: [unique, not_null]
--       - name: status
--         tests:
--           - accepted_values:
--               values: ['NEW', 'COMPLETE', 'CANCELLED']
--       - name: customer_id
--         tests:
--           - relationships:
--               to: ref('dim_customers')
--               field: customer_id
\end{lstlisting}

\textbf{Data Metric Functions} complement dbt's build-time tests with platform-native, always-on monitoring: attach DMFs (freshness, row count, null count, uniqueness) to tables and schedule evaluation---\texttt{ALTER TABLE ... SET DATA\_METRIC\_SCHEDULE = '60 MINUTE'}---with results queryable in \texttt{SNOWFLAKE.LOCAL.DATA\_METRIC\_FUNCTION\_RESULTS} \cite{snowflakeDmfDocs}. dbt tests gate \emph{deploys}; DMFs watch \emph{reality} between deploys.

\begin{tipbox}
dbt incremental models without \texttt{unique\_key} silently duplicate rows on re-runs. Declare the key, and add the \texttt{unique} test so CI catches what the config missed.
\end{tipbox}

\section{Wednesday: Migrations, Alerts, and Cost-Anomaly Detection}
\index{schemachange}\index{expand-contract}\index{cost anomaly}
\subsection{Versioned Migrations with schemachange}
Snowflake DDL has no built-in migration ledger. \textbf{schemachange} provides it: versioned scripts (\texttt{V1.3.2\_\_add\_order\_channel.sql}) applied in order, exactly once, tracked in a change-history table; repeatables (\texttt{R\_\_...}) reapply on content change for views and procedures \cite{schemachangeDocs}. CI runs schemachange per environment; the ledger is the audit trail of what changed where and when.

\subsection{Expand--Migrate--Contract for Zero Downtime}
Renaming a column the naive way---drop, add---breaks every consumer between the two statements. The professional sequence spans releases:
\begin{enumerate}
    \item \textbf{Expand:} add the new nullable column (\texttt{order\_total\_v2}) in one release.
    \item \textbf{Migrate:} dual-write and backfill; repoint dbt models and dashboards; keep a compatibility view (\texttt{CREATE OR REPLACE VIEW ... AS SELECT ..., order\_total\_v2 AS order\_total}) aliasing old names.
    \item \textbf{Contract:} in a later release, after a CI gate proves no dbt model or BI extract references the old column, drop it.
\end{enumerate}
In dbt, \texttt{on\_schema\_change} on incremental models participates: \texttt{append\_new\_columns} for the expand phase, \texttt{sync\_all\_columns} only after consumers migrate.

\begin{pitfall}
\textbf{Failure story:} a schemachange script dropped \texttt{order\_total} in the same release dbt introduced \texttt{order\_total\_v2}; the contract step ran before dbt rebuilt the shim view, and the executive revenue dashboard errored for 40 minutes. Rule: expand and contract in separate releases, keep the compatibility view until every consumer migrates, and gate the contract migration in CI.
\end{pitfall}

\subsection{Snowflake Alerts and Cost-Anomaly Detection}
\index{Snowflake Alerts}
An \textbf{alert} is a scheduled condition: every interval, evaluate a query; if it returns rows, execute an action---typically a notification-integration call to email or SNS-to-Slack \cite{snowflakeAlertsDocs}. The canonical first alert detects credit anomalies against \texttt{WAREHOUSE\_METERING\_HISTORY}:

\begin{lstlisting}[language=SQL, caption={Cost-anomaly alert: 2x deviation from trailing 7-day average}]
CREATE OR REPLACE ALERT gov.credit_anomaly_alert
  WAREHOUSE = gov_wh
  SCHEDULE = '60 MINUTE'
  IF (EXISTS (
    WITH daily AS (
      SELECT warehouse_name, DATE(start_time) AS d,
             SUM(credits_used) AS credits
      FROM SNOWFLAKE.ACCOUNT_USAGE.WAREHOUSE_METERING_HISTORY
      WHERE start_time >= DATEADD(day, -8, CURRENT_TIMESTAMP())
      GROUP BY 1, 2)
    SELECT d.warehouse_name, d.credits,
           AVG(p.credits) AS trailing_avg
    FROM daily d
    JOIN daily p
      ON p.warehouse_name = d.warehouse_name
     AND p.d BETWEEN DATEADD(day, -7, d.d) AND DATEADD(day, -1, d.d)
    WHERE d.d = CURRENT_DATE()
    GROUP BY 1, 2
    HAVING d.credits > 2 * AVG(p.credits)))
  THEN CALL SYSTEM$SEND_EMAIL(...);
\end{lstlisting}

\section{Thursday Hands-On Project: The CI/CD Pipeline with Observability}
\index{hands-on project!CI/CD}
This lab assembles the week: a GitHub repository where pull requests deploy migrations and test models, with cost observability wired end to end.

\subsection*{Scenario}
The platform team requires every Snowflake change reviewed, tested, and traceable: no manual DDL in production, failing data tests block deploys, and cost anomalies page within an hour.

\subsection*{Procedure}
\begin{enumerate}
    \item Create the repository: \texttt{migrations/} (schemachange), \texttt{models/} (dbt), Terraform for warehouses and monitors.
    \item Configure key-pair secrets in GitHub Actions (never passwords in CI variables).
    \item Build the workflow: on PR, run schemachange against QA, then \texttt{dbt build} (models + tests).
    \item Create the resource monitor with 75\%-notify / 100\%-suspend triggers in Terraform.
    \item Deploy the cost-anomaly alert from Wednesday wired to an email notification integration.
    \item Attach two DMFs (freshness, null count) to \texttt{fct\_orders} with a 60-minute schedule.
    \item Demonstrate failure: open a PR whose model breaks \texttt{not\_null}; show CI red and the deploy blocked.
    \item Walk an expand--contract column rename through two separate PRs with the shim view.
\end{enumerate}

\subsection*{Deliverables}
\begin{itemize}
    \item \textbf{Repository:} migrations, models with tests, Terraform, and the Actions workflow.
    \item \textbf{Evidence pack:} a green deploy run, the blocked bad-PR run, monitor configuration, and a fired test alert.
    \item \textbf{Migration ledger:} schemachange history table output for both environments.
    \item \textbf{Expand--contract record:} both PRs, the shim view, and the contract gate's consumer-reference check.
\end{itemize}

\subsection*{Acceptance Criteria}
\begin{itemize}
    \item Production DDL changes only through merged PRs; the ledger proves it.
    \item A failing dbt test blocks deploy; a failing migration never half-applies.
    \item The monitor and the anomaly alert are demonstrated live.
    \item The column rename completes across two releases with zero consumer-visible downtime.
\end{itemize}

\subsection*{Stretch Goals}
\begin{itemize}
    \item Add the contract-gate script failing CI while any model or extract references the old column.
    \item Extend the alert suite: DMF breach alerts and warehouse-queue-depth anomalies.
    \item Add a terraform plan comment bot to PRs for reviewer visibility.
\end{itemize}

\section{Friday Real-World Architecture: The Multi-Environment DataOps Framework}
\index{Real-World Architecture!DataOps}
A regulated enterprise runs Dev, QA, and Prod Snowflake accounts with change-freeze windows, audited promotion, and developer self-service. The framework's shape:

\begin{figure}[H]
    \centering
    \resizebox{\textwidth}{!}{%
    \begin{tikzpicture}[
        repo/.style={rectangle, rounded corners, draw=codegray, thick, fill=white, align=center, minimum width=2.5cm, minimum height=1.0cm, font=\small\bfseries},
        ci/.style={rectangle, rounded corners, draw=primary, thick, fill=primary!6, align=center, minimum width=2.5cm, minimum height=1.0cm, font=\small\bfseries},
        env/.style={rectangle, rounded corners, draw=codegreen!60!black, thick, fill=green!8, align=center, minimum width=2.2cm, minimum height=1.0cm, font=\small\bfseries},
        obs/.style={rectangle, rounded corners, draw=red!60!black, thick, fill=red!8, align=center, minimum width=2.5cm, minimum height=1.0cm, font=\small\bfseries},
        arrow/.style={-{Stealth[length=2.5mm]}, draw=primary, thick},
        lbl/.style={font=\scriptsize\itshape, text=codegray}
    ]
        \node[repo] (git) at (0, 0) {Git\\dbt + Terraform\\+ migrations};
        \node[ci] (pr) at (3.6, 1.2) {PR pipeline\\schemachange$\rightarrow$QA\\dbt build + tests};
        \node[ci] (rel) at (3.6, -1.2) {Release pipeline\\promote on tag\\contract gate};
        \node[env] (qa) at (7.4, 1.2) {QA};
        \node[env] (prod) at (7.4, -1.2) {Prod\\(freeze-aware)};
        \node[obs] (obs) at (11.2, 0) {Monitors + Alerts\\+ DMFs $\rightarrow$ Slack};
        \draw[arrow] (git) -- (pr);
        \draw[arrow] (git) -- (rel);
        \draw[arrow] (pr) -- (qa);
        \draw[arrow] (rel) -- (prod);
        \draw[arrow, dashed] (qa) |- (obs);
        \draw[arrow, dashed] (prod) |- (obs);
    \end{tikzpicture}%
    }
    \caption{Multi-environment DataOps: PRs prove changes in QA; tagged releases promote to Prod through a contract gate; observability wraps everything.}
    \label{fig:chapter23-dataops}
\end{figure}

\subsection{Operating Principles}
\begin{enumerate}
    \item \textbf{Everything is code.} Objects in Terraform, schema evolution in schemachange, transformations and tests in dbt; Snowsight edits in Prod are policy violations detected by drift reports.
    \item \textbf{Environments differ by data and blast radius, not shape.} Same Terraform modules, different variables: QA monitors suspend, Prod monitors page.
    \item \textbf{Zero-downtime is a process property.} Expand--contract with shim views and CI contract gates, rehearsed on every rename---not a hero move during incidents.
    \item \textbf{Observability closes the loop.} Resource monitors bound spend, DMFs watch data reality, alerts route anomalies to humans, and \texttt{ACCOUNT\_USAGE} retains the evidence.
\end{enumerate}

\section{Interview Q\&A}
\begin{enumerate}
    \item \textbf{Q: Why manage Snowflake objects with Terraform instead of UI clicks?}\\
    A: Declarative state in git gives review (plan), audit (apply logs), repeatability across environments, and drift detection; click-ops has none of these.
    \item \textbf{Q: Which four dbt generic tests should every model carry?}\\
    A: \texttt{unique}, \texttt{not\_null}, \texttt{accepted\_values}, and \texttt{relationships}.
    \item \textbf{Q: What does a Data Metric Function monitor?}\\
    A: Table-attached quality metrics---freshness, volume, null counts, uniqueness---evaluated on a schedule with results in \texttt{DATA\_METRIC\_FUNCTION\_RESULTS}.
    \item \textbf{Q: How does schemachange apply migrations safely?}\\
    A: Versioned scripts apply in order exactly once per environment, tracked in a change-history table; repeatables reapply only on content change.
    \item \textbf{Q: Walk through the expand-contract sequence for renaming a column.}\\
    A: Add the new column in release one; dual-write, backfill, and migrate consumers behind a compatibility view; drop the old column in a later release after CI proves no consumer references it.
\end{enumerate}

\section{Further Reading}
Terraform's Snowflake provider, dbt Core, and schemachange are the toolchain references \cite{terraformDocs,terraformSnowflakeProvider,dbtCoreDocs,schemachangeDocs}. Snowflake's resource-monitor, alert, and DMF documentation covers the platform-native observability \cite{snowflakeResourceMonitorsDocs,snowflakeAlertsDocs,snowflakeDmfDocs}. Chapter 10's audit patterns supply the queries these alerts automate.

\section*{Key Takeaways}
\begin{itemize}
    \item Everything is code: Terraform for objects, schemachange for evolution, dbt for transformation and tests.
    \item Resource monitors are cost circuit breakers; their posture differs per environment by design.
    \item dbt tests gate deploys; DMFs watch reality between deploys.
    \item Expand--contract with shim views makes schema change boring; one-release renames are outages.
    \item Alerts over \texttt{ACCOUNT\_USAGE} turn cost and quality anomalies into pages within an hour.
    \item CI secrets use key-pair auth; no passwords in CI variables, ever.
\end{itemize}

\section{Exercises}
\begin{enumerate}
    \item \textbf{(Conceptual)} Explain the difference between build-time tests (dbt) and runtime monitors (DMFs).
    \item \textbf{(Conceptual)} Why must expand and contract be separate releases even when both scripts are ready?
    \item \textbf{(Hands-on)} Add a singular dbt test enforcing a business invariant (e.g., refunds never exceed order totals) and break it deliberately.
    \item \textbf{(Hands-on)} Write the drift report comparing Terraform state against \texttt{SHOW} output for warehouses.
    \item \textbf{(Hands-on)} Build the freeze-window gate: releases blocked outside approved windows except hotfix labels.
    \item \textbf{(Design)} Design the environment-promotion model for a team with feature branches and a shared QA.
    \item \textbf{(Design)} Specify the notification taxonomy: which alerts page, which post to Slack, which email a digest.
    \item \textbf{(Operations)} Write the monthly DataOps health report: deploy frequency, failure rate, mean recovery, alert precision.
    \item \textbf{(Cost)} Model credit savings from aggressive dev-environment auto-suspend and monitor suspension.
    \item \textbf{(Interview)} Argue to a skeptical manager why ``slow'' expand--contract deploys are cheaper than fast breaking ones.
\end{enumerate}
